\section{纯 $\SU(2)$ 理论的 Seiberg-Witten 解}\label{sec:pureSU2}\index{PureSU2@纯 $\SU(2)$ 理论|hyperpage}
我们终于做好了充分准备，可以开始分析最简单的非阿贝尔 $\cN{=}2$ 超对称理论，即纯 $\SU(2)$ 规范理论。我们主要遵循原始论文 \cite{Seiberg:1994rs} 的论述，不过将采用文献 \cite{Martinec:1995by} 首次给出的 \SeibergWitten\ 曲线形式，因为它更适合后面的推广。

\subsection{单圈跑动与无穷远处的单值变换}\index{Monodromy@单值变换|hyperpage}\index{UPlane@$u$ 平面|hyperpage}
纯 $\SU(2)$ 理论只包含规范群 $\SU(2)$ 的一个 $\cN{=}2$ 矢量多重态，其拉格朗日量由式 \eqref{vectorlag} 给出。为便于查阅，我们在此将其重写如下：
\begin{equation}
L=\frac{\Im\tau}{4\pi} \int d^4\theta \tr \Phi^\dagger e^{[V,\cdot]} \Phi +\int d^2\theta \frac{-\ii}{8\pi}\tau \tr W_\alpha W^\alpha + cc. \label{high}
\end{equation}

在经典层面，超对称真空由 D 项约束的解来刻画：
\begin{equation}
[\Phi^\dagger,\Phi]=0.
\end{equation}
这意味着 $\Phi$ 可以通过规范旋转对角化。令
\begin{equation}
\Phi=\diag(a,-a).\label{dv}
\end{equation}

粗略地说，规范耦合 $\tau$ 按照 $\SU(2)$ 理论的单圈重整化，从一个很高的能标跑动到能标 $a$。随后，真空期望值 $a$ 将规范群 $\SU(2)$ 破缺为 $\U(1)$。存在带未破缺 $\U(1)$ 荷的有质量激发，但它们很快就会退耦，因此在能标 $a$ 以下，耦合几乎保持不变。这一演化如图 \ref{fig:roughrunning} 所示。
\begin{figure}[h]
\[
\inc{su2pure-runningX}
\]
\caption{耦合跑动的示意图。\label{fig:roughrunning}}
\end{figure}

让我们稍微更定量地描述这一过程。我们对 $\U(1)$ 拉格朗日量和规范耦合的归一化由式 \eqref{u1coupling} 和 \eqref{u1tau} 给出，在此重写如下：
\begin{equation}
\frac{1}{2e^2} F_{\mu\nu}^{\U(1)}F_{\mu\nu}^{\U(1)} + \frac{\theta}{16\pi^2} F_{\mu\nu}^{\U(1)} \tilde F_{\mu\nu}^{\U(1)},
\qquad \text{以及}\qquad
\tau^{\U(1)}= \frac{4\pi i}{e^2} + \frac{\theta}{2\pi}.
\end{equation}
在破缺真空中，低能的 $\U(1)$ 与高能的 $\SU(2)$ 之间具有式 \eqref{u1embed} 所示的关系，我们也将其重写如下：
\begin{equation}
F_{\mu\nu}^{\SU(2)}=\diag(F_{\mu\nu}^{\U(1)},-F_{\mu\nu}^{\U(1)}).
\end{equation}
将其代入高能拉格朗日量 \eqref{high}，并比较各个 $\tau$ 的定义，我们得到
\begin{equation}
\tau^{\U(1)} = 2\tau^{\SU(2)}.
\end{equation}
量子修正会改变这一关系。

当真空期望值由式 \eqref{dv} 给出时，以 $\tau(a)$ 表示 $\U(1)$ 规范场的低能耦合，并以 $\tau_{UV}$ 表示高能重整化点 $\Lambda_{UV}$ 处 $\SU(2)$ 规范场的高能耦合。单圈跑动 \eqref{n2running} 随即给出
\begin{align}
\tau(a)&=2\tau_{UV}-\frac{8}{2\pi i}\log \frac{a}{\Lambda_{UV}}+\cdots \\
&=-\frac{8}{2\pi i}\log \frac{a}{\Lambda} + \cdots\label{puretau}
\end{align}
其中我们定义了
\begin{equation}
\Lambda^4 = \Lambda_{UV}^4 e^{2\pi i \tau_{UV}}.\label{lambdaSU2pure}
\end{equation}
将式 \eqref{puretau} 积分一次即可得到对偶变量 $a_D$，结果为
\begin{equation}
a_D=-\frac{8a}{2\pi i}\log\frac{a}{\Lambda} + \cdots.
\end{equation}
只要保持 $|a|\gg |\Lambda|$，耦合 $\tau(a)$ 就仍然很弱，上述计算便给出了可靠的近似。

一种以规范不变方式标记超对称真空的方法是使用
\begin{equation}
u=\frac12\vev{\tr\phi^2}=a^2+\cdots
\end{equation}
其中 $\cdots$ 表示量子修正。考虑将 $u$ 的相位绝热地旋转 $2\pi$：
\begin{equation}
u=e^{i\theta}|u|, \qquad \theta=0 \sim 2\pi
\end{equation}
于是有 $a\mapsto -a$。根据 $a_D$ 的显式形式，可得 $a_D\to -a_D+4a$。我们将其记作
\begin{equation}
(a,a_D)\to (a,a_D) \begin{pmatrix}
-1 & 4 \\
0 & -1
\end{pmatrix}.\label{amat}
\end{equation}
BPS 粒子的质量公式为
\begin{equation}
M=|na+ma_D|=\left|
(a,a_D)
\begin{pmatrix}
n \\ m
\end{pmatrix}
\right|.
\end{equation}
因此，变换 \eqref{m_inf_nf_0} 也可以归结为荷的变换：
\begin{equation}
\begin{pmatrix}
n \\ m
\end{pmatrix}
\to
\begin{pmatrix}
-1 & 4 \\
0 & -1
\end{pmatrix}
\begin{pmatrix}
n \\ m
\end{pmatrix}.\label{chargemat}
\end{equation}

我们将矩阵
\begin{equation}
M_\infty= \begin{pmatrix}
-1 & 4 \\
0 & -1
\end{pmatrix}\label{m_inf_nf_0}
\end{equation}
称为无穷远处的单值变换。图 \ref{fig:m_inf_nf_0} 示意性地描绘了这一情形。由 $u$ 参数化的超对称真空空间通常称为 $u$ 平面。
\begin{figure}[h]
\[
\inc{su2pure-monodromy-fromafar}
\]
\caption{无穷远处的单值变换。\label{fig:m_inf_nf_0}}
\end{figure}

在我们的论证中，矩阵 \eqref{amat} 原本可能包含非整数矩阵元，因为这些矩阵元是从 $a$ 和 $a_D$ 的近似公式中读出的。然而，变换 \eqref{chargemat} 必须将整数向量映为整数向量，这就保证了矩阵 \eqref{chargemat} 是整数矩阵。不仅如此，这一变换只是在重新标记荷，因此不应改变狄拉克配对
\begin{equation}
nm'-mn'=\det \begin{pmatrix}
n & n' \\
m & m'
\end{pmatrix}
\end{equation}
当存在两个荷分别为 $(n,m)$ 和 $(n',m')$ 的粒子时，它衡量空间中所携带的角动量。形如
\begin{equation}
\begin{pmatrix}
n \\ m
\end{pmatrix}
\to
M
\begin{pmatrix}
n \\ m
\end{pmatrix}
\end{equation}
的变换对狄拉克配对的影响为
\begin{equation}
\det \begin{pmatrix}
n & n' \\
m & m'
\end{pmatrix}
\to
\det M
\det
\begin{pmatrix}
n & n' \\
m & m'
\end{pmatrix}.
\end{equation}
因此，$M$ 的行列式必须为一。行列式为一的 $2\times 2$ 整数矩阵称为 $\SL(2,\bZ)$ 的一个元素。矩阵 \eqref{m_inf_nf_0} 确实满足这一条件，令人安心。

\subsection{强耦合区域中的行为}
让我们研究 $u$ 平面内部的强耦合区域中发生了什么。要实现全纯函数 $a$ 和 $a_D$ 的单值变换 $M_\infty$，这一内部区域中至少需要有一个奇点。因而，最朴素的预期便是图 \ref{fig:m_0_nf_0} 所示的结构。
\begin{figure}[h]
\[
\inc{su2pure-monodromy-naive}
\]
\caption{一个行不通的朴素猜测。\label{fig:m_0_nf_0}}
\end{figure}
奇点会在哪里？这里，系统中未破缺的离散 $\U(1)$ R 对称性可以派上用场。回忆一下，我们的 $\cN{=}2$ 理论具有 $\SU(2)_R$ 对称性。在经典层面，还可以考虑一个 $\U(1)_R$ 对称性，其标准 R 荷赋值如下：
\begin{equation}
\begin{array}{r|ccc}
R=0 & & A\\
1& \lambda & & \lambda \\
2 & & \Phi
\end{array}.
\end{equation}
同一超对称多重态中的不同分量具有不同的荷，因此这是一个 R 对称性。

在量子层面，旋转
\begin{equation}
\lambda\to e^{i\varphi} \lambda,
\end{equation}
具有反常，但正如第 \ref{sec:anomaly} 小节所述，可以通过
\begin{equation}
\theta_{UV}\to \theta_{UV}+8\varphi,
\end{equation}
予以补偿。

因此，$\varphi=\pi/4$ 是一个真正的对称性，其作用为
\begin{equation}
\theta\to \theta+2\pi,\quad
\Phi\to e^{\pi i/2}\Phi.
\end{equation}
它生成系统的 $\bZ_4$ 离散 R 对称性。对低能变量，其作用为
\begin{equation}
\theta_{IR}\to \theta_{IR}+4\pi,\qquad
u\to -u.\label{discrete_r_nf_0}
\end{equation}
那么，如果 $u=u_0$ 处有一个奇点，$u=-u_0$ 处就应当还有一个奇点。因此，如果只有一个奇点，它就位于 $u=0$。

如果确实如此，我们会发现 $\tau(a)$ 由下式给出：
\begin{equation}
\tau(a)=-\frac{8}{2\pi i}\log \frac{a}{\Lambda} + f(a).\label{wrongt}
\end{equation}
其中 $f(a)$ 是一个唯一奇点位于 $a=0$ 的亚纯函数。然而，这听起来并不正确。耦合由 $\Im\tau(a)$ 给出，它是一个全纯函数的虚部。因而，它没有下界，所以会在某些 $a$ 值处变成负数。这意味着耦合 $g^2$ 在那些位置为负，系统变得不稳定。例如，假设 $f(a)=0$，当 $|a|$ 足够小时，虚部就是负的。我们由此得出结论：假设只在 $u=0$ 处有一个奇点，实在过于朴素。

其次简单的可能性是，假设 $u=\pm u_0$ 处有两个奇点，见图 \ref{fig:monodromy_nf_0_correct}。
\begin{figure}[h]
\[
\inc{su2pure-monodromy-correct}
\]
\caption{接下来的猜测，事实证明是正确的。\label{fig:monodromy_nf_0_correct}}
\end{figure}
系统中唯一的能标是动力学能标 $\Lambda$，因此 $u_0$ 应当由 $c\Lambda$ 给出，其中 $c$ 是一个数。将绕两个奇点的单值变换记为 $M_\pm$，则应有
\begin{equation}
M_\infty=M_+ M_-,
\end{equation}
因为绕 $u$ 平面无穷远处的路径，在拓扑上等同于先绕 $u=-u_0$、再绕 $u=u_0$ 的路径。由于两个奇点通过一个对称性彼此交换，绕它们的单值变换本质上应当相同，只差荷的重新标记。等价地说，它们应当通过某个 $\SL(2,\bZ)$ 矩阵 $X$ 共轭：
\begin{equation}
M_-=X M_+ X^{-1}
\end{equation}
注意，矩阵 $X$ 可以看作与对称操作 \eqref{discrete_r_nf_0} 对应的半单值变换，见图 \ref{fig:discrete_r_nf_0}。
\begin{figure}[h]
\[
\inc{mplus-mminus}
\]
\caption{$M_+$ 与 $M_-$ 之间的关系。\label{fig:discrete_r_nf_0}}
\end{figure}

这些方程的一个解为
\begin{equation}
M_+=STS^{-1}=\begin{pmatrix}
1 & 0\\
-1 & 1
\end{pmatrix}, \quad
M_-=T^2 S T S^{-1} T^{-2}=\begin{pmatrix}
-1 & 4 \\
-1 & 3
\end{pmatrix}\label{mpm}
\end{equation}
其中 $S$ 和 $T$ 由式 \eqref{S}、\eqref{T} 给出。注意，我们有
\begin{equation}
X=T^2,
\end{equation}
这与离散 R 对称性 \eqref{discrete_r_nf_0} 将 $\theta_{IR}$ 平移 $4\pi$ 这一事实大体相容。

\subsection{Seiberg-Witten 解}\label{sec:purecurve}\index{SeibergWittenCurve@Seiberg-Witten 曲线|hyperpage}\index{SeibergWittenDifferential@Seiberg-Witten 微分|hyperpage}\index{UltravioletCurve@紫外曲线|hyperpage}\index{Period@周期|hyperpage}\index{ComplexStructure@复结构|hyperpage}
让我们显式构造满足这些单值变换的全纯函数 $a$ 和 $a_D$。注意，一个全纯函数由它的奇点唯一确定。因而，假定我们正确地识别了奇点，只要找到一个在定义域的奇点附近和无穷远处具有正确性质的候选函数，它就必然是正确答案本身。因此，只需构造一个候选函数，再检验它是否满足条件即可。

\subsubsection{曲线}
我们首先引入两个辅助复变量 $x$ 和 $z$，然后考虑方程
\begin{equation}
\Sigma:\quad
\Lambda^2 z + \frac{\Lambda^2}z=x^2-u.\label{curve_pure}
\end{equation}
我们将此方程看作在由 $x$ 和 $z$ 构成的复二维空间中定义了一个复一维子空间。\footnote{以 $(z,x)$ 作为坐标遵循了文献 \cite{Martinec:1995by} 的约定。将我们所称的 $(z,x)$ 写为 $(t,v)$ 也很常见，这一记法源自文献 \cite{Witten:1997sc}。} 随着 $u$ 改变，方程也随之改变，因此这个子空间的形状也会变化。这一复一维对象称为 \SeibergWitten\ 曲线。\footnote{它是实二维的，因此从通常的观点来看是一个曲面。数学家很奇怪，他们将一维对象都视为曲线，无论它是复一维还是实一维。} 一个称为 Seiberg-Witten 微分的微分形式
\begin{equation}
\lambda=x\frac{dz}z,
\end{equation}
将在后面发挥重要作用。\footnote{到目前为止，符号 $\lambda$ 一直表示伴随表示中的费米子；不过从现在起，除非另有说明，我们主要用 $\lambda$ 表示微分形式。}

由 $z$ 参数化的空间本身也很重要。我们在 $z$ 复平面上加入 $z=\infty$ 这一点；等价地说，将 $z$ 视为球面的复坐标。将此球面记为 $C$，并称其为这一系统的\Gaiotto 曲线。作为 $z$ 的函数，变量 $x$ 有四个平方根分支点，见图 \ref{fig:gaiotto_curve_nf_0}。
\begin{figure}[h]
\[
\inc{su2purecurveX}
\]
\caption{纯 $\SU(2)$ 理论的\Gaiotto 曲线 $C$。\label{fig:gaiotto_curve_nf_0}}
\end{figure}
\begin{figure}[h]
\[
\inc{su2purecurve-sheets}
\]
\caption{纯 $\SU(2)$ 理论的 \SeibergWitten\ 曲线 $\Sigma$ 的各个叶片。\label{fig:curve_nf_0}}
\end{figure}

于是，曲线 $\Sigma$ 是 $C$ 的二重覆盖，
\begin{equation}
\Sigma\stackrel{2:1}{\longrightarrow} C,
\end{equation}
见图 \ref{fig:curve_nf_0}。随后，我们如图所示在曲线上画出两个一维闭链 $A$、$B$，并规定
\begin{equation}
a=\frac{1}{2\pi i}\oint_A \lambda,\qquad
a_D=\frac{1}{2\pi i}\oint_B \lambda.
\end{equation}

让我们检验由此定义的函数 $a(u)$ 和 $a_D(u)$ 是否满足物理上预期的性质。首先，计算 $\tau(a)$：
\begin{equation}
\tau(a)=\frac{\partial a_D}{\partial a}
=\frac{\partial a_D/\partial u}{\partial a/\partial u}.
\end{equation}
对 $u$ 的导数可以如下计算：
\begin{align}
\frac{\partial a}{\partial u}&=\int_A \frac{\partial}{\partial u} \lambda
=\int_A \frac{dz}{xz}\\
\frac{\partial a_D}{\partial u}&=\int_B \frac{\partial}{\partial u} \lambda
=\int_B \frac{dz}{xz}
\end{align}
其中积分号内对 $u$ 求导时固定 $z$。微分形式 $\omega=dz/(xz)$ 在 $\Sigma$ 上处处有限，即使在看似危险的点 $z=0$、$z=\infty$ 或 $x=0$ 处也是如此。例如，当 $x=0$ 时，有 $z\sim c+c'x^2$，其中 $c$ 和 $c'$ 为常数。于是 $dz/(xz)\sim (2c'/c)dx$。
\begin{figure}[h]
\[
\inc{su2purecurve-torus}\qquad \Longrightarrow \qquad
\inc{su2purecurve-parallelogramX}
\]
\caption{纯 $\SU(2)$ 理论的 \SeibergWitten\ 曲线 $\Sigma$ 经光滑化后是一个环面。\label{fig:para_nf_0}}
\end{figure}

给定曲线 $\Sigma$ 上一条从固定点 $P_0$ 出发的开路径，就得到一个从路径终点到另一个复平面的映射：
\begin{equation}
t=\int_{P_0}^P \omega.\label{mapping_to_t}
\end{equation}
如图 \ref{fig:para_nf_0} 所示，曲线 $\Sigma$ 被映到复平面中的一个平行四边形，它的边界是闭链 $A$ 和 $B$ 的像所构成的线段。像式 \eqref{mapping_to_t} 这样的全纯映射都保持角度。因此，闭链 $B$ 的像始终位于闭链 $A$ 的像的左侧。于是
\begin{equation}
\tau(a)=\frac{\partial a_D/\partial u}{\partial a/\partial u}=
\frac{\int_B {dz}/(xz)}{\int_A {dz}/(xz)}
\end{equation}
的取值位于实轴的左侧，因此
\begin{equation}
\Im \tau(a)>0,
\end{equation}
这就保证耦合的平方 $g^2(a)$ 始终为正。复数 $\tau(a)$ 称为环面的周期或复结构。

\subsubsection{单值变换 $M_\infty$}
让我们检验曲线 \eqref{curve_pure} 是否重现了从物理考虑确定的单值变换。将曲线 $\Sigma$ 写为
\begin{equation}
z+\frac1z=\frac{x^2}{\Lambda^2}-\frac{u}{\Lambda^2}.\label{simplified_eq}
\end{equation}
由此可见，当 $|u|\gg\Lambda^2$ 时，函数 $x(z)$ 的两个分支点 $z_\pm$ 大致位于
\begin{equation}
z_+\sim -u/\Lambda^2,\qquad z_-\sim -\Lambda^2/u.
\end{equation}
此外，$z=0$ 和 $z=\infty$ 也是分支点。我们选取的分支割线分别从 $z=0$ 延伸到 $z=z_-$，以及从 $z=z_+$ 延伸到 $z=\infty$。

将 $A$ 闭链放在 $|z|=1$ 处。沿它的积分非常容易计算：在 $|z|=1$ 附近，$x\simeq\sqrt{u}$，因此
\begin{equation}
a=\frac{1}{2\pi i}\oint x\frac{dz}z\simeq \sqrt{u}.
\end{equation}
至于沿 $B$ 闭链的积分，主要贡献来自变量 $z$ 不太靠近分支点的区域。变量 $x$ 同样可以近似为 $\sqrt{u}\simeq a$，因此
\begin{equation}
a_D=\frac{2}{2\pi i} \int^{z_-}_{z_+} x\frac{dz}z\simeq \frac{2\cdot 2}{2\pi i}  \int^1_{u/\Lambda^2} a\frac{dz}z \simeq -\frac{8a}{2\pi i}\log \frac{a}{\Lambda}.\label{aDpure}
\end{equation}
由这两个方程可知，通过曲线 $\Sigma$ 定义的 $a$ 和 $a_D$ 在 $u\sim\infty$ 周围具有正确的单值变换，
\begin{equation}
M_\infty=\begin{pmatrix}
-1 & 4 \\
0  & -1
\end{pmatrix}.
\end{equation}

通过更仔细的计算，可以显式求出式 \eqref{aDpure} 或其导数 $\tau(a)$ 的修正。从曲线 \eqref{curve_pure} 的形式可以明显看出，修正可以按 $\Lambda^2$ 的幂展开，但实际上它们由 $\Lambda^4$ 的幂给出。我们得到
\begin{equation}
\tau(a)= -\frac{8}{2\pi i}\log\frac{a}{\Lambda} + \sum_{k=0}^\infty c_k \left(\frac{\Lambda}{a}\right)^{4k}
\label{tauexpansion}
\end{equation}
其中 $c_k$ 是无量纲有理数。现在，我们知道了式 \eqref{puretau} 中以 $\cdots$ 隐去的各项。例如，可以通过引入 $\tilde z=\Lambda^2 z$ 来理解这一展开。此时曲线为 $\tilde z+\Lambda^4/\tilde z=x^2-u$；将 $\Lambda^4/\tilde z$ 视为对曲线领头阶形式 $\tilde z=x^2-u$ 的微扰，就可以计算 $a$ 和 $\tilde a_D$。

让我们解释这些按 $\Lambda^4$ 的幂给出的修正。从式 \eqref{lambdaSU2pure} 可知，$\Lambda^{4k}$ 项携带相位 $e^{ik\theta_{UV}}$，其中 $\theta_{UV}$ 是 $\theta$ 角。正如式 \eqref{instanton_contri} 所述，它对应于瞬子数为 $k$ 的构型。这一展开显式说明，低能耦合 $\tau(a)$ 唯一的微扰修正来自单圈阶，同时还存在来自瞬子的非微扰修正。在瞬子背景下进行真正的路径积分计算，应当重现系数 $c_k$。对于单瞬子贡献 $c_1$，文献 \cite{Finnell:1995dr} 已经完成了这一计算。随后，文献 \cite{Nekrasov:2002qd,Nekrasov:2003rj} 将其推广到了所有 $k$。总之，我们看到，各种物理量都由单圈对数贡献加上瞬子修正的组合给出。得益于始于文献 \cite{Nekrasov:2002qd} 的一系列发展，现在已知它们在瞬子展开的所有阶都相符。在附录 \ref{sec:instcount} 中，我们将直接由曲线计算系数 $c_k$，并看到它们与微观瞬子计算的结果一致。

\subsubsection{单值变换 $M_\pm$}\index{MonopolePoint@单极子点|hyperpage}\index{DyonPoint@双荷子点|hyperpage}
接下来研究强耦合奇点周围的单值变换。再看一下式 \eqref{simplified_eq}，显然，当 $z+1/z=\pm2$ 时，情形颇为特殊。当 $u=2\Lambda^2$ 时，两个分支点在 $z_\pm=-1$ 处碰撞；当 $u=-2\Lambda^2$ 时，它们在 $z_\pm=+1$ 处碰撞。这就是图 \ref{fig:monodromy_nf_0_correct} 中引入的奇点 $u=\pm u_0$。

以 $u=u_0=2\Lambda^2$ 附近的行为为例。令 $u=2\Lambda^2+\delta u$，则分支点位于
\begin{equation}
z_\pm-1\propto \pm \sqrt{\delta u}.
\end{equation}
分支点 $z_\pm$ 与闭链 $A$、$B$ 的局部放大图见图 \ref{fig:closeup_nf_0}。
\begin{figure}[h]
\[
\inc{monodromy-on-cyclesX}
\]
\caption{单极子点周围的单值变换对闭链的作用。\label{fig:closeup_nf_0}}
\end{figure}
当我们绕 $u=2\Lambda^2$ 缓慢改变 $u$ 的值时，两个分支点 $z=z_\pm$ 发生交换。这使闭链 $A$ 如图所示地改变，等价于原来的闭链 $A$ 减去闭链 $B$。闭链 $B$ 显然不变。因此有
\begin{equation}
a \to a-a_D,\qquad
a_D\to a_D
\end{equation}
等价地，单值变换为
\begin{equation}
M_+=\begin{pmatrix}
1 & 0 \\
-1 & 1
\end{pmatrix},
\end{equation}
重现了式 \eqref{mpm}。

让我们研究 $u=u_0=2\Lambda^2$ 处的物理。施行 $S$ 变换 \eqref{S}，
\begin{equation}
a'=-a_D,\quad
a_D'=a
\end{equation}
以交换电荷和磁荷。它们作为 $u$ 的函数可写为
\begin{equation}
a'=c(u-u_0),\quad
a_D'=\frac{a'}{2\pi i}\log c'(u-u_0)\label{cc}
\end{equation}
其中 $c$ 和 $c'$ 是两个常数，由此可得
\begin{equation}
\tau_D(a')=\frac{\partial a_D'}{\partial a'}\sim +\frac{\log a'}{2\pi i}.\label{dual_run}
\end{equation}
注意，$a'$ 确定了系统的能标。这一结果与含有一个带荷超多重态的 $\cN{=}2$ 超对称 $\U(1)$ 规范理论的耦合跑动具有相同的行为；这个超多重态由 $\cN{=}1$ 手征多重态 $(Q,\tilde Q)$ 组成。此时的超势耦合为
\begin{equation}
\int d^2\theta Qa'\tilde Q.
\end{equation}
写成
\begin{equation}
\tau_D=\frac{4\pi i}{g_D^2} + \frac{\theta_D}{2\pi},
\end{equation}
就可发现，当 $u\to u_0$ 时，
\begin{equation}
g_D\to 0
\end{equation}
成立。$Q$ 的量子的质量由 BPS 质量公式给出：
\begin{equation}
\text{$Q$ 的量子的质量}=|a'|=|a_D|.
\end{equation}
因此，我们将带荷手征多重态 $Q$ 识别为原理论中单极子的二次量子化版本。在弱耦合区域中非常重的单极子，现在变得非常轻。

施行离散 R 对称性 \eqref{discrete_r_nf_0}，就很容易得到 $u=-u_0$ 处的行为。由于映射由 $T^2$ 给出，我们发现此时非常轻的粒子具有电荷 $n=2$ 和磁荷 $m=1$，也就是说，它们是双荷子。出于这些原因，$u=u_0=2\Lambda$ 这一点通常称为单极子点，而 $u=-u_0=-2\Lambda$ 这一点称为双荷子点。
% [Tutor Note: 原文此处印作 $u=u_0=2\Lambda$、$u=-u_0=-2\Lambda$，量纲与本节前文不符：本小节已确定强耦合奇点位于 $u=\pm 2\Lambda^2$（见 "When $u=2\Lambda^2$, the two branch points collide at $z_\pm=-1$"），故单极子点、双荷子点应为 $u=\pm 2\Lambda^2$。原文 $2\Lambda$ 系笔误，正文照原文排，读者请以 $2\Lambda^2$ 为准。]

\subsection{超对称性较少的情形}
在继续研究 $\cN{=}2$ 系统之前，让我们稍作停顿，看看能从纯 $\cN{=}2$ $\SU(2)$ 理论的解中，了解到超对称性较少的理论的哪些性质。本小节考虑的一般拉格朗日量为
\begin{equation}
\begin{aligned}
&(\int d^2\theta \frac{-i}{8\pi}\tau \tr W_\alpha W^\alpha + cc.)
+\frac{\Im \tau}{4\pi}\int d^4\theta \Phi^\dagger \Phi \\
&\qquad +(\int d^2\theta \frac{m}2\tr\Phi^2 + cc.)
+(\mu \lambda_\alpha\lambda^\alpha + cc.)
\end{aligned}\label{441}
\end{equation}
当 $m=\mu=0$ 时，这一系统具有 $\cN{=}2$ 超对称性。令 $|m|\to\infty$，手征超场 $\Phi$ 退耦，最终得到第 \ref{sec:n=1pure} 小节讨论过的 $\cN{=}1$ 纯 $\SU(2)$ 理论。接着令 $|\mu|\to\infty$，规范微子 $\lambda$ 退耦，就恢复了纯玻色的杨--米尔斯理论。

\subsubsection{$\cN{=}1$ 系统}\index{MonopoleCondensation@单极子凝聚|hyperpage}\index{Confinement@禁闭|hyperpage}
首先考虑 $\cN{=}1$ 系统。当 $m$ 很小时，$m\tr\Phi^2$ 项可以视为对刚刚得到的 $\cN{=}2$ 解的微扰。用变量 $u$ 表示时，$\int d^2\theta m\tr\Phi^2$ 项为 $\sim\int d^2\theta mu$，因此，除非 $u$ 位于奇点，对 $u$ 的 F 项方程就不可能满足。在一般的 $u$ 值处，没有超对称真空。

当 $u$ 接近 $u_0=2\Lambda^2$ 时，超势中存在额外项
\begin{equation}
\int d^2\theta Qa'\tilde Q = \int d^2\theta c(u-u_0)Q\tilde Q
\end{equation}
其中常数 $c$ 已在式 \eqref{cc} 中引入。加上 $\int d^2\theta mu$ 项，对 $u$、$Q$ 和 $\tilde Q$ 的 F 项方程分别为
\begin{equation}
m=cQ\tilde Q,\quad
(u-u_0)\tilde Q=0,\quad
(u-u_0)Q=0.
\end{equation}
于是得到一个解
\begin{equation}
u=u_0,\quad Q\tilde Q=m/c.
\end{equation}
真空被固定在 $u=u_0$，并且存在非零的单极子凝聚 $Q\tilde Q=m/c$。在 $u=-u_0$ 处作类似的论证可知，还存在另一个超对称真空，
\begin{equation}
u=-u_0,\quad Q'\tilde Q'=m/c
\end{equation}
其中 $Q'$、$\tilde Q'$ 是双荷子场。

总之，我们在 $u=\pm u_0$ 处找到了两个超对称真空，在其中单极子或双荷子发生凝聚，从而具体实现了禁闭由带磁荷对象的凝聚所导致这一思想，见图 \ref{fig:less_supersymmetric}。
\begin{figure}[h]
\[
\inc{vacuaX}
\]
\caption{软破缺的 $\cN{=}1$、$\cN{=}0$ 理论的真空。\label{fig:less_supersymmetric}}
\end{figure}

回忆一下，由反常破缺的连续 R 对称性
\begin{equation}
\Phi\to e^{i\varphi}\Phi,
\end{equation}
可以通过
\begin{equation}
\theta_{UV}\to\theta_{UV}+4\varphi.
\end{equation}
予以补偿。将它应用到拉格朗日量 \eqref{441}，可以看到
\begin{equation}
m\vev{\tr\Phi^2}=\frac{-i}{2\pi}\vev{\tr W^\alpha W_\alpha}
\end{equation}
由此得到
\begin{equation}
\vev{\lambda_\alpha\lambda^\alpha}\propto\pm2\pi i m\Lambda^2=: \pm\Lambda_{\cN{=}1}^3.
\end{equation}
必须记住，右端包含相位 $e^{i\theta_{UV}/2}$。

现在保持 $\Lambda_{\cN{=}1}$ 不变，并取极限 $m\to\infty$。这应当给出纯 $\cN{=}1$ $\SU(2)$ 杨--米尔斯理论。令人安心的是，与第 \ref{sec:n=1pure} 小节一样，我们在这里也看到了两个真空。

\subsubsection{纯玻色系统}\index{SoftBreaking@软破缺|hyperpage}\index{FirstOrderPhaseTransition@一级相变|hyperpage}
现在令 $\mu\neq0$，同时保持 $|\mu|\ll|\Lambda_{\cN{=}1}|$。在这一极限下，规范微子质量项 $\mu\lambda_\alpha\lambda^\alpha$ 的效应由一阶微扰论给出，真空能为
\begin{equation}
V\propto\Re(\pm\mu\Lambda^3_{\cN{=}1})
\propto\Lambda^4_{\cN{=}0}\Re(\pm e^{i\theta_{UV}/2}).
\end{equation}
文献 \cite{Konishi:1996iz} 首次指出了这一点。

我们看到，$\cN{=}1$ 超对称理论的两个简并真空分裂为两个具有不同能量密度的能级，分别对应单极子凝聚与双荷子凝聚。将 $\theta_{UV}$ 从 $0$ 缓慢改变到 $2\pi$，会交换这两个能级，它们在 $\theta_{UV}=\pi$ 处交叉。因而，至少在 $|\mu|$ 足够小时，$\theta_{UV}=\pi$ 处存在一级相变。

一个有趣的问题是，在 $|\mu|\to\infty$ 的极限下，也就是在纯玻色的杨--米尔斯理论中，这一一级相变是否仍然存在。让我们给出支持其持续存在的论证。思路是将两个带磁荷或双荷的外部粒子之间势能的行为用作序参量 \cite{Aharony:2013hda}。

首先更仔细地考虑动力学。两个分支的区别在于发生凝聚的粒子种类不同，因此可相应地称它们为单极子支和双荷子支。按照我们的约定，这些粒子的荷分别为 $(n,m)=(0,1)$ 和 $(2,1)$。在我们的归一化下，$\SU(2)$ 伴随场对于未破缺 $\U(1)$ 对称性的荷为 $(2,0)$。由于不存在荷为 $(1,0)$ 的动力学粒子，单极子的荷 $(0,1)$ 是最小允许磁荷的两倍。因此，这种外部单极子的荷可以写为 $(n,m)=(0,1/2)$。

首先考虑引入两个荷为 $(n,m)=(1,0)$ 的外部电粒子。在两个分支中，凝聚的单极子或双荷子都会将电场束缚成通量管。通量管具有恒定张力，并且不能成对产生动力学粒子，因为所有动力学粒子的荷都是 $(\pm2,0)$。因此，通量管不会断裂，势能是线性的。荷为 $(1,0)$ 的电粒子被禁闭。

换一种情形，考虑在系统中引入外部单极子，并测量两者之间的势能。在 $\theta=0$ 时，不失一般性，可以假设单极子支的能量更低。背景中有荷为 $(n,m)=(0,1)$ 的动力学单极子粒子发生凝聚。现在引入两个荷为 $(n,m)=(0,1/2)$ 的外部单极子。荷为 $(n,m)=(0,1)$ 的外部粒子所产生的磁场被屏蔽并呈指数衰减。因而，它们之间的势能基本上是常数。
% [Tutor Note: 原文此句 "The magnetic field produced by the external particles with charges $(n,m)=(0,1)$" 中的荷 $(0,1)$ 与前一句刚引入的外部单极子的荷 $(0,1/2)$ 不一致：此处被屏蔽的应是外部荷 $(0,1/2)$ 的磁场，而 $(0,1)$ 是背景中凝聚的动力学单极子的荷；疑原文应作 $(0,1/2)$（或径指动力学粒子）。正文照原文排，特此提示。]

再考虑向单极子支引入两个荷为 $(1,1/2)$ 的外部粒子。动力学单极子无法屏蔽电荷，因此电场被禁闭在通量管内。两者之间的势能是线性的，它们受到禁闭。

可以在双荷子支重复这一分析。外部粒子之间势能的行为可以总结如下：
\begin{center}
\begin{tabular}{r|cc}
& $(0,1/2)$ & $(1,1/2)$ \\
\hline
单极子支 & 屏蔽 & 禁闭 \\
双荷子支 & 禁闭 & 屏蔽
\end{tabular}
\end{center}
当 $\theta$ 从 $0$ 缓慢且连续地改变到 $2\pi$ 时，这两种行为相互交换。因此，至少应当存在一次相变。通过格点强耦合展开或计算机模拟来验证这一分析，将是很有趣的。

\subsection{$\SU(2)$ 与 $\SO(3)$ 的比较}\index{GlobalGaugeGroup@规范群的整体形式|hyperpage}
此时，讨论一个有关理论的 \SeibergWitten\ 曲线的微妙问题或许会有所帮助：它依赖于究竟将规范群选为 $\SU(2)$ 还是 $\SO(3)$。初次阅读时可以跳过这一小节。
% [Tutor Note: 原文 "a rather sutble point" 中 "sutble" 为 "subtle" 的拼写笔误（同一段下文即出现 "subtle manner"），译文按 "微妙" 译出。]

在本节中，我们对荷的选择是：
\begin{equation}
(n,m)=(1,0)
\end{equation}
表示 $\SU(2)$ 中的一个电二重态，
\begin{equation}
(n,m)=(0,1)
\end{equation}
表示与 $\SU(2)\to\U(1)$ 破缺相联系的一个 't Hooft-Polyakov 单极子。

不过，理论中所有动力学粒子的荷都具有形式
\begin{equation}
(n,m)=(2k,m)\label{dync}
\end{equation}
其中 $k$ 和 $m$ 为整数。此外，由于没有任何处于 $\SU(2)$ 二重态中的动力学场，我们可以考虑一个荷为
\begin{equation}
(n,m)=(0,1/2).
\end{equation}
的外部单极子。相对于理论中的任意动力学粒子，它仍然满足狄拉克量子化条件；这些动力学粒子的荷由式 \eqref{dync} 给出。
% [Tutor Note: 原文 "whose chages are given by \eqref{dync}" 中 "chages" 为 "charges" 的拼写笔误，译文按 "荷" 译出。]

相应地，单值变换矩阵
\begin{equation}
M_\infty=\begin{pmatrix}
-1 & 4 \\
0 & -1
\end{pmatrix},\qquad
M_+=\begin{pmatrix}
1 & 0 \\
-1 & 1
\end{pmatrix},\qquad
M_-=\begin{pmatrix}
-1 & 4 \\
-1 & 3
\end{pmatrix}
\end{equation}
的右上角矩阵元都是 $4$ 的整数倍。

因此，可以作如下处理。通过
\begin{equation}
(n',m')=(n/2,2m)
\end{equation}
定义重新缩放后的电荷与磁荷，单值变换仍然取整数值：
\begin{equation}
M_\infty'=\begin{pmatrix}
-1 & 1 \\
0 & -1
\end{pmatrix},\qquad
M_+'=\begin{pmatrix}
1 & 0 \\
-4 & 1
\end{pmatrix},\qquad
M_-'=\begin{pmatrix}
-1 & 1 \\
-4 & 3
\end{pmatrix}.
\end{equation}
现在，BPS 质量公式为
\begin{equation}
|na+ma_D|=|n'a'+m'a'_D|\qquad\text{其中}\quad(a',a'_D)=(2a,a_D/2)
\end{equation}
相应地，新的 $A$、$B$ 闭链与旧闭链之间的关系为
\begin{equation}
A'=2A,\qquad B'=B/2.
\end{equation}
与图 \ref{fig:para_nf_0} 一样，将相应的 \SeibergWitten\ 曲线视为复平面的商，便得到图 \ref{fig:so3}。
\begin{figure}[h]
\[
\inc{so3}
\]
\caption{\SeibergWitten\ 曲线的不同选择。从左到右依次为：纯 $\SU(2)$ 理论、纯 $\SO(3)$ 理论，以及文献 \cite{Seiberg:1994rs} 中的选择。\label{fig:so3}}
\end{figure}

标准解释是，将闭链为 $A$、$B$ 的 \SeibergWitten\ 曲线视为纯 $\SU(2)$ 理论的曲线，而将闭链为 $A'=2A$、$B'=B/2$ 的 \SeibergWitten\ 曲线视为纯 $\SO(3)$ 理论的曲线。二者的区别以一种相当微妙的方式表现出来。

对于规范群 $\SO(3)$，单极子点处的单值变换为上面给出的 $M_+'$。这意味着该处的轻粒子对于低能 $\U(1)$ 场的荷为 $2$：左下角的矩阵元 $-4$ 由荷的平方给出。其原因在于，低能 $\U(1)$ 嵌入的是 $\SO(3)$ 而不是 $\SU(2)$，因此其周期缩短了一半。单极子的荷为 $(n',m')=(0,2)$，而 $n'$ 与 $m'$ 的最大公因数为二。

在双荷子点，单值变换 $M_-'$ 仍与 $\begin{pmatrix}1&1\\0&1\end{pmatrix}$ 共轭，也就是说，轻双荷子对于 $\U(1)$ 的荷为 $1$。双荷子的荷为 $(n',m')=(1,2)$，$n'$ 与 $m'$ 的最大公因数为一。因此，单极子点与双荷子点在物理上的等价性不再成立。当系统被置于 $\bR^3\times S^1$ 或更复杂的流形上时，这些微妙区别会对系统产生更显著的影响。

另一个有趣的事实是，$M_\infty'$、$M_+'$ 和 $M_-'$ 这一组合，与含有 $3$ 个无质量味的 $\SU(2)$ 理论的单值变换矩阵完全相同，后者将在第 \ref{sec:nf3} 小节讨论。然而，正如后面将看到的，纯 $\SO(3)$ 理论与含有 $3$ 个无质量味的 $\SU(2)$ 理论，其物理却截然不同。这说明了一个显而易见的道理：仅凭库仑支的结构，并不能确定整个理论。

最后需要提及，在关于纯 $\SU(2)$ 理论的最初论文 \cite{Seiberg:1994rs} 中，使用的是另一条以 $2A$ 和 $B$ 为两个闭链的曲线，如图 \ref{fig:so3} 所示。这一选择适合动力学粒子的谱 \eqref{dync}，但现在已知，当考虑非平凡流形上的理论以及线算符的性质时，它并没有很充分的理由。因此，最好采用纯 $\SU(2)$ 曲线或纯 $\SO(3)$ 曲线，即图 \ref{fig:so3} 中的前两项。
% [Tutor Note: 原文 "in the very orignal paper" 中 "orignal" 为 "original" 的拼写笔误，译文按 "最初论文" 译出。]
